\documentclass[10pt,a4paper]{article}
\usepackage[margin=2.2cm]{geometry}
\usepackage{fontspec}
\setmonofont{DejaVu Sans Mono}[Scale=0.84]
\newfontfamily\unifont{DejaVu Serif}[Scale=0.9]
\usepackage{amsmath,amssymb,mathtools}
\setcounter{MaxMatrixCols}{20}
\usepackage{longtable,booktabs,array}
\usepackage{fancyvrb}
\usepackage[dvipsnames]{xcolor}
\usepackage[colorlinks=true,linkcolor=NavyBlue,urlcolor=NavyBlue]{hyperref}
\usepackage{microtype}
\usepackage{needspace}
\usepackage{tocloft}\setlength{\cftbeforesecskip}{0.45em}
\setlength{\parindent}{0pt}\setlength{\parskip}{0.55em}
\setlength{\emergencystretch}{8em}
\makeatletter
\renewcommand\section{\@startsection{section}{1}{\z@}{-3.5ex \@plus -1ex \@minus -.2ex}{2.3ex \@plus.2ex}{\normalfont\Large\bfseries\raggedright}}
\renewcommand\subsection{\@startsection{subsection}{2}{\z@}{-3.25ex\@plus -1ex \@minus -.2ex}{1.5ex \@plus .2ex}{\normalfont\large\bfseries\raggedright}}
\renewcommand\subsubsection{\@startsection{subsubsection}{3}{\z@}{-3.25ex\@plus -1ex \@minus -.2ex}{1.5ex \@plus .2ex}{\normalfont\normalsize\bfseries\raggedright}}
\makeatother
\renewcommand{\arraystretch}{1.15}
\allowdisplaybreaks
\begin{document}
\title{dirac16complex as dark matter and dark energy: symbols and their definitions\\[0.4em]\large A dictionary of every symbol used in the analysis of section 24 of the notebook 7+1.nb, which rests on the solutions of the coupled field equations in the three gravity regimes}
\author{Repository: github.com/once-ere/7plus1 — prepared for Nash\\ Notebook: 7+1.nb (Wolfram Language 15.0.1)\\ Document generated by tools/make\_docs.py on 2026-10-28}
\date{}
\maketitle
\begin{abstract}
For each symbol this dictionary gives the mathematical notation, the name used in the Mathematica notebook 7+1.nb (where there is one; the names of section 24 begin with \texttt{dk}), the defining equation and its meaning, grouped by topic. Longer defining equations follow each table as numbered displayed formulas. Every equation needed to define a symbol is written in this document.
\end{abstract}
\tableofcontents
\Needspace{12\baselineskip}\section*{1. Conventions}\addcontentsline{toc}{section}{1. Conventions}
\begin{itemize}
\item Units: $\hbar=c=1$; energies, masses and temperatures in GeV unless stated; cross sections in cm$^2$, lifetimes in s; $\hbar c=1.973269804\times10^{-14}$ GeV cm, $\hbar=6.582119569\times10^{-25}$ GeV s, 1 GeV $=1.78266192\times10^{-24}$ g. The uniform Kohn–Sham matter of section 8 below uses units $m=1$.
\item Coordinates $X=(x_0,x_1,\dots,x_7)$: $x_0$ is time (the direction of the octonion unit $e_0=1$), $x_1,x_2,x_3$ ordinary space, $x_4,\dots,x_7$ the four auxiliary space-like coordinates; on the world brane $x=(x_0,x_1,x_2,x_3)$, and $y=(x_4,\dots,x_7)$.
\item A prime is $d/dx_0$; $\tau$ is the time since the singularity; $t$ the cosmic time of the brane; "ln" the natural logarithm.
\item Status labels: DERIVED (an exact consequence of the field equations, checked), ESTIMATE (order of magnitude, assumptions stated), OPEN (not decided by the model).
\end{itemize}
\textbf{The signature: $\{7,1,0\}$ and $\mathrm{diag}(+1,-1,\dots,-1)$.} The 7+1 spacetime has one time-like dimension ($x_0$, the direction of the octonion unit $e_0=1$) and seven space-like dimensions ($x_1,x_2,x_3$ ordinary space, $x_4,\dots,x_7$ auxiliary space). Two different objects describe this, and neither may be replaced by the other. (i) The tensor package xAct, used in section 12 of 7+1.nb exactly as in the author’s notebook Einstein-Lovelock-7+1-Nash.nb, declares the metric with \texttt{DefMetric[\{7, 1, 0\}, g71[-a, -b], CD, …]}: the list $\{7,1,0\}$ declares seven space-like dimensions, one time-like dimension and no degenerate dimension — the numbers of entries $+1$ and $-1$ of $\eta_{71}=\mathrm{diag}(-1,+1,\dots,+1)$ (\texttt{η71}), the same spacetime written in the sign convention of 4+4.nb (space-like $=+1$, time-like $=-1$). (ii) The matrix $\eta=-\eta_{71}=\mathrm{diag}(+1,-1,\dots,-1)$ (\texttt{η17}) is simply the squares of the octonion basis: $e_0^2=+1$ for the unit $e_0=1$ (time) and $e_i^2=-1$ for the seven imaginary units $e_1,\dots,e_7$ (space); more generally it is the real part of the product of two basis units, and it gives the quadratic form of the octonions:
\begin{gather*}
\eta_{AB}=\mathrm{Re}(e_Ae_B),\qquad e_0^2=+1,\qquad e_i^2=-1\ \ (i=1,\dots,7),\\ \mathrm{Re}(x\,x)=x_0^2-x_1^2-\cdots-x_7^2=\eta_{AB}\,x^Ax^B\qquad (x=x^Ae_A\in\mathbb O).
\end{gather*}
$\eta$ gives the signs of the Clifford relation $\Gamma_A\Gamma_B+\Gamma_B\Gamma_A=2\,\eta_{AB}\,1_{16}$ of the real 16×16 Dirac matrices: time squares to $+1_{16}$ and each of the seven space directions to $-1_{16}$, exactly like the octonion units (in the notation $Cl(p,q)$, with $p$ generators squaring to $+1$ and $q$ to $-1$, the algebra $Cl(1,7)\cong\mathbb R(16)$). It also gives the signs of the metric components: $g_{00}>0$ along the time-like direction, negative squares along the space-like directions. Every component computed by 7+1.nb follows from the component matrix of the metric; the declared list fixes for xAct the sign of the determinant of the metric, \texttt{SignDetOfMetric[g71]} $=-1$, which is the same in both sign conventions, since $\det \eta=\det \eta_{71}=-1$. 7+1.nb checks \texttt{SignatureOfMetric[g71]} $=\{7,1,0\}$, \texttt{SignDetOfMetric[g71]} $=-1$, $\mathrm{Re}(e_Ae_B)=\eta_{AB}$ and $\mathrm{Re}(x\,x)=\eta_{AB}x^Ax^B$.

\Needspace{12\baselineskip}\section*{2. The model}\addcontentsline{toc}{section}{2. The model}
{\small
\begin{longtable}{>{\raggedright\arraybackslash}p{\dimexpr 0.1700\linewidth-2\tabcolsep\relax}>{\raggedright\arraybackslash}p{\dimexpr 0.2100\linewidth-2\tabcolsep\relax}>{\raggedright\arraybackslash}p{\dimexpr 0.3400\linewidth-2\tabcolsep\relax}>{\raggedright\arraybackslash}p{\dimexpr 0.2800\linewidth-2\tabcolsep\relax}}
\toprule
\textbf{symbol} & \textbf{Mathematica} & \textbf{definition} & \textbf{meaning} \\
\midrule\endfirsthead
\toprule
\textbf{symbol} & \textbf{Mathematica} & \textbf{definition} & \textbf{meaning} \\
\midrule\endhead
\endfoot
\bottomrule\endlastfoot
$\psi$, $\bar\psi$ & \texttt{dirac16complex} & $\bar\psi=\psi^\dagger\Gamma_0$ & the 16-component complex field and its Dirac conjugate \\
$\Gamma_A$, $\Gamma^A$, $\Gamma^\mu$ & \texttt{Γ17[A]}, \texttt{Γ17up[A]} & (2.1) & real 16×16 Dirac matrices; frame matrices \\
$\Omega_\mu$ & — & (2.1) & spin connection \\
$m$ & \texttt{mDM} (symbolic, section 24.2) &  & mass of dirac16complex (free) \\
$g_r$, $\Theta^K$, $\eta_K$ & \texttt{gg[r]} & (2.2) & the nine self-couplings; Clifford products; their signs \\
$\alpha_k$, $\Lambda$, $\kappa$ & \texttt{alf[k]}, \texttt{lamb}, \texttt{kapS} & $\alpha_0=2\Lambda$, $\alpha_1=1$, $\kappa=8\pi G_8=M_8^{-6}$ & Lovelock couplings, bulk cosmological constant \\
$L_2$, $L_3$ & — & (2.3) & Gauss–Bonnet and cubic Lovelock densities \\
$U(b)$, $U_0$, $K_U$, $b_0$ & \texttt{uPot}, \texttt{u0B}, \texttt{kU} & (2.4) & stabilizing potential of the auxiliary space \\
$W=\lbrace y=0\rbrace$, $h$ & — & $h_{\mu\nu}=g_{\mu\nu}(x,0)$ & world brane and its induced metric \\
$f_W$ & \texttt{fWb} &  & brane tension \\
$\mathcal L_{\rm SM}$, $\Phi$, $\tilde\Phi$ & — & $\tilde\Phi=i\sigma_2\Phi^*$ & Standard Model Lagrangian on the brane, Higgs doublet \\
$y_{is}$, $G_P^f$, $\lambda_P$ & — & (2.5) & Yukawa (neutrino), vector and Higgs portals \\
$L_i$, $\psi_s$, $f$ & — &  & lepton doublets; brane-right-handed parts of $\psi$ at the brane; Standard Model fermions \\*
\end{longtable}}
\begin{align*}
\Gamma_A\Gamma_B+\Gamma_B\Gamma_A&=2\eta_{AB}1_{16},\quad \Gamma^A=\eta^{AA}\Gamma_A,\quad \Gamma^\mu=e_A{}^\mu\Gamma^A,&&(2.1) \\
\Omega_\mu&=\tfrac14\eta_{AC}\omega_\mu{}^C{}_B\Gamma^A\Gamma^B,\quad g_{\mu\nu}=\eta_{AB}e^A{}_\mu e^B{}_\nu, \\
\Theta^K&=i^{\sigma_r}\Gamma^{A_1}\cdots\Gamma^{A_r},\quad K=(A_1<\dots<A_r),&&(2.2) \\
\sigma_r&=\tfrac{r(r-1)}2\bmod2,\quad \eta_K=\eta_{A_1A_1}\cdots\eta_{A_rA_r}, \\
L_2&=R^2-4R_{\mu\nu}R^{\mu\nu}+R_{\mu\nu\rho\sigma}R^{\mu\nu\rho\sigma},&&(2.3) \\
L_3&=\tfrac18\delta^{\alpha_1\beta_1\alpha_2\beta_2\alpha_3\beta_3}_{\gamma_1\delta_1\gamma_2\delta_2\gamma_3\delta_3}R^{\gamma_1\delta_1}{}_{\alpha_1\beta_1}R^{\gamma_2\delta_2}{}_{\alpha_2\beta_2}R^{\gamma_3\delta_3}{}_{\alpha_3\beta_3}, \\
U(b)&=U_0+\tfrac{K_U}2\ln^2(b/b_0),\qquad b^8=\det(-g_{kl}),\ k,l=4,\dots,7,&&(2.4) \\
S_P&=\int_Wd^4x\sqrt{|h|}\Big[-\sum_{i,s}y_{is}\bar L_i\tilde\Phi P_R\psi_s+\text{h.c.}&&(2.5) \\
&\qquad-G_P^f(\bar\psi\Gamma^\mu\psi)(\bar f\gamma_\mu f)-\lambda_P\Phi^\dagger\Phi\,\bar\psi\psi\Big].
\end{align*}
\Needspace{12\baselineskip}\section*{3. Cosmology of the three regimes}\addcontentsline{toc}{section}{3. Cosmology of the three regimes}
{\small
\begin{longtable}{>{\raggedright\arraybackslash}p{\dimexpr 0.1700\linewidth-2\tabcolsep\relax}>{\raggedright\arraybackslash}p{\dimexpr 0.2100\linewidth-2\tabcolsep\relax}>{\raggedright\arraybackslash}p{\dimexpr 0.3400\linewidth-2\tabcolsep\relax}>{\raggedright\arraybackslash}p{\dimexpr 0.2800\linewidth-2\tabcolsep\relax}}
\toprule
\textbf{symbol} & \textbf{Mathematica} & \textbf{definition} & \textbf{meaning} \\
\midrule\endfirsthead
\toprule
\textbf{symbol} & \textbf{Mathematica} & \textbf{definition} & \textbf{meaning} \\
\midrule\endhead
\endfoot
\bottomrule\endlastfoot
$a$, $b$ & \texttt{a3}, \texttt{b4}; \texttt{aS}, \texttt{bS} (symbols) & (3.1) & scale factors of ordinary and auxiliary space \\
$H_a$, $H_b$ & \texttt{Ha}, \texttt{Hb} & $a'/a$, $b'/b$ & Hubble rates \\
$\beta$, $m_r$, $\phi$ & \texttt{betR}, \texttt{rIIm} & $\ln(b/b_0)$; $m_r^2=\kappa K_U/12$; $\sqrt{12}M_{\rm Pl}\beta$ & radion, its mass, canonical radion \\
$\varepsilon$, $p_a$, $p_b$ & \texttt{feStress[W]} & (3.2) & energy density and pressures in ordinary and auxiliary space \\
$W(a,b)$ & \texttt{wW} & (3.2) & comoving energy of a source at fixed occupation numbers \\
$\omega$, $k_a$, $k_b$ & \texttt{kaP}, \texttt{kbP} & (3.3) & comoving energy and momenta of one particle \\
$V_\perp$, $G_4$, $\Lambda_4$, $M_{\rm Pl}$ & \texttt{rIIVp}, \texttt{rIIMPl} & (3.4) & volume of the auxiliary torus; four-dimensional constants \\
$\rho_4$, $\rho_{\rm DM}$ & — & $\rho_4=V_\perp\varepsilon$ & energy per brane volume; that of the dark matter \\
$\delta\beta$ & \texttt{dkDbeta} & (3.5) & radion displacement by the dark matter today \\
$h_*$, $p$, $w_a$, $w_b$ & \texttt{rIhStar} & (3.6) & matter-dominated cubic regime; equation-of-state ratios \\
$P$, $Q$, $\mathcal C$ & \texttt{rIfl} & (3.7) & Einstein power laws with a fluid \\
$A$, $\vartheta$, $\tau_E$, $H_E$ & \texttt{rIIaF}, \texttt{rIIphF}, \texttt{rIItEF} & $\beta\approx A\sin(m_r\tau+\vartheta)/(m_r(\tau-\tau_E))$, $H_E=H_a+2H_b$ & radion oscillation after capture \\
$f$ & \texttt{dkF} $=$ \texttt{rIIfrVals[[2]]} &  & share of the dirac16complex gas in the energy after capture \\
$H_0$, $\Omega_r$, $\Omega_m$, $\Omega_b$, $\Omega_c$, $\Omega_\Lambda$ & \texttt{rIIH0g}, \texttt{rIIOr}, \texttt{rIIOm}, \texttt{dkOc}, \texttt{rIIOL} & (3.8) & Hubble constant and density fractions today \\
$h$ & \texttt{dkh} & $H_0=100h$ km/s/Mpc & reduced Hubble constant (0.674) \\
$H_\infty$ & \texttt{rIIIH} & $H_0\Omega_\Lambda^{1/2}$ & final de Sitter rate \\
$t_0$ & \texttt{dkAge} & 13.79 Gyr & age of the universe \\*
\end{longtable}}
\begin{align*}
ds^2&=dx_0^2-a^2(dx_1^2+dx_2^2+dx_3^2)-b^2(dx_4^2+\cdots+dx_7^2),&&(3.1) \\
\varepsilon&=\frac W{a^3b^4},\qquad p_a=-\frac{\partial_aW}{3a^2b^4},\qquad p_b=-\frac{\partial_bW}{4a^3b^3},\qquad \varepsilon^{\prime}+3H_a(\varepsilon+p_a)+4H_b(\varepsilon+p_b)=0,&&(3.2) \\
\omega&=\sqrt{m^2+k_a^2/a^2+k_b^2/b^2},&&(3.3) \\
V_\perp&=b_0^4L_c^4,\qquad 8\pi G_4=\frac\kappa{V_\perp}=\frac1{M_{\rm Pl}^2},\qquad \Lambda_4=\Lambda+\kappa U_0+\frac{\kappa f_W}{V_\perp},\qquad 3H_a^2=8\pi G_4\rho_4+\Lambda_4,&&(3.4) \\
\beta^{\prime\prime}&+3H_a\beta^{\prime}+4\beta^{\prime2}=\kappa\Big(\frac\varepsilon6-\frac{p_a}2+\frac{p_b}3\Big)+\frac\Lambda3,\qquad \delta\beta=\frac{\kappa\varepsilon}{6m_r^2}=\frac{H_0^2\Omega_c}{2m_r^2},&&(3.5) \\
h_*&=-\sqrt{\alpha_2/(54\alpha_3)},\qquad a\propto\tau^p,\quad p=\frac1{1+w_a},\qquad p_a=w_a\varepsilon,\quad p_b=w_b\varepsilon,&&(3.6) \\
3P(1+w_a)&+4Q(1+w_b)=2,\qquad 3(P^2+4PQ+2Q^2)=\mathcal C,\qquad a\propto\tau^P,\ b\propto\tau^Q,&&(3.7) \\
H_a&=H_0\sqrt{\Omega_ra^{-4}+\Omega_ma^{-3}+\Omega_\Lambda},\qquad \Omega_m=\Omega_b+\Omega_c,\qquad \Omega_c=\frac{\Omega_ch^2}{h^2}.&&(3.8)
\end{align*}
\Needspace{12\baselineskip}\section*{4. Kaluza–Klein modes of the auxiliary torus}\addcontentsline{toc}{section}{4. Kaluza–Klein modes of the auxiliary torus}
{\small
\begin{longtable}{>{\raggedright\arraybackslash}p{\dimexpr 0.1700\linewidth-2\tabcolsep\relax}>{\raggedright\arraybackslash}p{\dimexpr 0.2100\linewidth-2\tabcolsep\relax}>{\raggedright\arraybackslash}p{\dimexpr 0.3400\linewidth-2\tabcolsep\relax}>{\raggedright\arraybackslash}p{\dimexpr 0.2800\linewidth-2\tabcolsep\relax}}
\toprule
\textbf{symbol} & \textbf{Mathematica} & \textbf{definition} & \textbf{meaning} \\
\midrule\endfirsthead
\toprule
\textbf{symbol} & \textbf{Mathematica} & \textbf{definition} & \textbf{meaning} \\
\midrule\endhead
\endfoot
\bottomrule\endlastfoot
$\vec n$, $\vec q$ & — & $\vec q=2\pi\vec n/L_c$, $\vec n\in\mathbb Z^4$ & Kaluza–Klein momentum \\
$\ell_\perp$ & \texttt{dkLp} & $V_\perp^{1/4}$ & size of the auxiliary torus \\
$\mu_n$, $M_n$ & \texttt{rIImuK} & (4.1) & Kaluza–Klein mass and mass of a level \\
$r_4(k)$, $\sigma(k)$ & \texttt{rIIr4}, \texttt{rIIsig} & (4.2) & number of $\vec n$ with $|\vec n|^2=k$; divisor sum \\
$N(\mu)$, $dN$ & — & (4.3) & number and density of levels \\
$\chi$, $\chi_n$ & — & $\psi(x,y)=\sum_n\chi_n(x)e^{i\vec q_n\cdot\vec y}/\sqrt{V_\perp}$ & four-dimensional fields of the levels (four brane Dirac spinors each); $\chi=\chi_0$ the lightest mode \\
$g$ & — & 16 & internal states of the lightest mode (four spinors × two spins × particle, antiparticle) \\*
\end{longtable}}
\begin{align*}
\mu_n&=\frac{2\pi|\vec n|}{\ell_\perp},\qquad M_n=\sqrt{m^2+\mu_n^2},&&(4.1) \\
r_4(k)&=8\sigma(k)-32\sigma(k/4),\qquad \sigma(k)=\sum_{d\mid k}d,\qquad \sigma(k/4)=0\ \text{unless}\ 4\mid k,&&(4.2) \\
N(\mu)&=\frac{V_\perp\mu^4}{32\pi^2},\qquad dN=\frac{V_\perp\mu^3d\mu}{8\pi^2}\qquad\text{(continuum limit)}.&&(4.3)
\end{align*}
\Needspace{12\baselineskip}\section*{5. Observational quantities}\addcontentsline{toc}{section}{5. Observational quantities}
{\small
\begin{longtable}{>{\raggedright\arraybackslash}p{\dimexpr 0.1700\linewidth-2\tabcolsep\relax}>{\raggedright\arraybackslash}p{\dimexpr 0.2100\linewidth-2\tabcolsep\relax}>{\raggedright\arraybackslash}p{\dimexpr 0.3400\linewidth-2\tabcolsep\relax}>{\raggedright\arraybackslash}p{\dimexpr 0.2800\linewidth-2\tabcolsep\relax}}
\toprule
\textbf{symbol} & \textbf{Mathematica} & \textbf{definition} & \textbf{meaning} \\
\midrule\endfirsthead
\toprule
\textbf{symbol} & \textbf{Mathematica} & \textbf{definition} & \textbf{meaning} \\
\midrule\endhead
\endfoot
\bottomrule\endlastfoot
$Y$ & — & $n/s$ & comoving number density (number per entropy) \\
$mY$ & \texttt{dkYm} & (5.1) & required mass per entropy of the dark matter \\
$\rho_c$, $s_0$, $g_{*s0}$, $T_{\rm CMB}$ & \texttt{dkRhoc}, \texttt{dks0}, \texttt{dkGs0}, \texttt{rIITcmb} & (5.1) & critical density, entropy density today and its degrees of freedom, CMB temperature \\
$M_D$, $R$, $n$ & \texttt{dkMD}; \texttt{rRS} & (5.2) & Planck scale and radius in the convention of the Particle Data Group \\
$\sigma_{\rm SI}$ & \texttt{dkSigSI[c\_\allowbreak{}P, m]} & (6.2) & spin-independent dark matter–nucleon cross section \\
$\sigma/m$ & \texttt{dkSigM[c\_\allowbreak{}r, m]} & (6.3) & self-interaction cross section per mass \\
$w_0$, $w_a$ & — & $w(a)=w_0+w_a(1-a)$ & parameters of evolving dark energy (DESI) \\
$m_{\rm WDM}$ & — &  & mass of a thermal-relic warm dark-matter particle (Lyman-$\alpha$ bound) \\
Pauli bound & \texttt{dkPauli} & $0.13\,\text{keV}\,(2/g)^{1/4}$ & phase-space bound on fermionic dark matter \\*
\end{longtable}}
\begin{align*}
mY&=\frac{\Omega_ch^2\,(\rho_c/h^2)}{s_0},\qquad \frac{\rho_c}{h^2}=3\Big(100\,\frac{\rm km}{\rm s\,Mpc}\Big)^2M_{\rm Pl}^2,\qquad s_0=\frac{2\pi^2}{45}g_{*s0}T_{\rm CMB}^3,\qquad g_{*s0}=2+\frac{21}4\cdot\frac4{11},&&(5.1) \\
M_{\rm Pl}^2&=R^nM_D^{n+2},\quad V_\perp=(2\pi R)^n\ (n=4)\quad\Rightarrow\quad M_D=(2\pi)^{2/3}M_8.&&(5.2)
\end{align*}
\Needspace{12\baselineskip}\section*{6. The lightest mode: couplings, detection, stability}\addcontentsline{toc}{section}{6. The lightest mode: couplings, detection, stability}
{\small
\begin{longtable}{>{\raggedright\arraybackslash}p{\dimexpr 0.1700\linewidth-2\tabcolsep\relax}>{\raggedright\arraybackslash}p{\dimexpr 0.2100\linewidth-2\tabcolsep\relax}>{\raggedright\arraybackslash}p{\dimexpr 0.3400\linewidth-2\tabcolsep\relax}>{\raggedright\arraybackslash}p{\dimexpr 0.2800\linewidth-2\tabcolsep\relax}}
\toprule
\textbf{symbol} & \textbf{Mathematica} & \textbf{definition} & \textbf{meaning} \\
\midrule\endfirsthead
\toprule
\textbf{symbol} & \textbf{Mathematica} & \textbf{definition} & \textbf{meaning} \\
\midrule\endhead
\endfoot
\bottomrule\endlastfoot
$c_P$, $c_y$, $c_V$, $c_r$ & \texttt{cPS}, \texttt{cyS}, \texttt{cVS}, \texttt{crS} & (6.1) & dimensionless couplings in units of $M_8$ (natural size: 1) \\
$\lambda_{\rm eff}$, $y_{\rm eff}$, $G_{\rm eff}$, $G_r$ & \texttt{dkLam}, \texttt{dkYeff}, —, \texttt{dkGr} & (6.1) & four-dimensional couplings of a level \\
$v$, $m_h$, $G_F$, $\alpha$ & \texttt{vEW}, \texttt{mHb}, \texttt{gFermi} & $v=(\sqrt2G_F)^{-1/2}$ & Higgs vacuum value, Higgs mass, Fermi constant, fine-structure constant \\
$f_N$, $m_N$, $\mu$ & — & 0.30, 0.9383 GeV, $mm_N/(m+m_N)$ & Higgs–nucleon coupling factor, nucleon mass, reduced mass \\
$G_{\chi N}$ & — & $\lambda_{\rm eff}f_Nm_N/m_h^2$ & contact coupling of $\chi$ to nucleons \\
$\mu_P$, $\theta$ & \texttt{dkMuP[c\_\allowbreak{}y]} & $\mu_P=y_{\rm eff}v/\sqrt2$, $\theta=\mu_P/m$ & portal Dirac mass; neutrino mixing angle \\
$\Gamma(\chi\to\nu\gamma)$, $\Gamma(\chi\to3\nu)$ & \texttt{dkTauG}, \texttt{dkTau3} (lifetimes) & (6.4) & decay rates through the Yukawa portal \\
$K$ & \texttt{dkK} & $\tau(\chi\to3\nu)\,c_y^2m^3$ & lifetime constant (s GeV$^3$) \\
$m_{\max}$ & \texttt{dkMmax} & $(K/(c_y^2t_0))^{1/3}$ & largest mass that outlives the universe ($c_y=1$) \\*
\end{longtable}}
\begin{align*}
\lambda_P&=\frac{c_P}{M_8^5},\quad y=\frac{c_y}{M_8^2},\quad G_P=\frac{c_V}{M_8^6},\quad g_r=\frac{c_r}{M_8^6};&&(6.1) \\
\lambda_{\rm eff}&=\frac{c_PM_8}{M_{\rm Pl}^2},\quad y_{\rm eff}=\frac{c_yM_8}{M_{\rm Pl}},\quad G_{\rm eff}=\frac{c_V}{M_{\rm Pl}^2},\quad G_r=\frac{c_r}{M_{\rm Pl}^2}, \\
\sigma_{\rm SI}&=\frac{\lambda_{\rm eff}^2f_N^2m_N^2\mu^2}{\pi m_h^4},&&(6.2) \\
\sigma&\approx\frac{G_r^2m^2}{4\pi},&&(6.3) \\
\Gamma(\chi\to\nu\gamma)&=\frac{9\alpha G_F^2m^5\sin^22\theta}{1024\pi^4},\qquad \Gamma(\chi\to3\nu)=\frac{G_F^2m^5\theta^2}{96\pi^3},&&(6.4) \\
\frac{\Gamma(\chi\to\nu\gamma)}{\Gamma(\chi\to3\nu)}&=\frac{27\alpha}{8\pi}.
\end{align*}
\Needspace{12\baselineskip}\section*{7. The abundance}\addcontentsline{toc}{section}{7. The abundance}
{\small
\begin{longtable}{>{\raggedright\arraybackslash}p{\dimexpr 0.1700\linewidth-2\tabcolsep\relax}>{\raggedright\arraybackslash}p{\dimexpr 0.2100\linewidth-2\tabcolsep\relax}>{\raggedright\arraybackslash}p{\dimexpr 0.3400\linewidth-2\tabcolsep\relax}>{\raggedright\arraybackslash}p{\dimexpr 0.2800\linewidth-2\tabcolsep\relax}}
\toprule
\textbf{symbol} & \textbf{Mathematica} & \textbf{definition} & \textbf{meaning} \\
\midrule\endfirsthead
\toprule
\textbf{symbol} & \textbf{Mathematica} & \textbf{definition} & \textbf{meaning} \\
\midrule\endhead
\endfoot
\bottomrule\endlastfoot
$O$ & \texttt{dkOver[T\_\allowbreak{}RH]} & (7.1) & overabundance of the primordial gas over the dark matter \\
$T_{\rm RH}$ & \texttt{trh} &  & reheating temperature of the brane \\
$g_*$, $g_{*s}$ & \texttt{rIIgstars} & $\rho_{\rm rad}=\frac{\pi^2}{30}g_*T^4$, $s=\frac{2\pi^2}{45}g_{*s}T^3$ & relativistic degrees of freedom \\
$\Gamma_{\rm brane}$, $\Gamma_{\rm bulk}$, $\Gamma_n$ & \texttt{rIIgam}, \texttt{rIIgBulk} & (7.2) & radion decay rates into brane fields and into Kaluza–Klein pairs \\
$b_{\rm QCD}$, $\alpha_s$, $N_c$ & — & $11-\frac23n_f$ & QCD beta-function coefficient, strong coupling, colours \\
$B(x,y)$ & \texttt{Beta} & $\Gamma(x)\Gamma(y)/\Gamma(x+y)$ & Euler beta function \\
$T_*$ & \texttt{dkTstar} & $(M_8^6/M_{\rm Pl})^{1/5}$ & normalcy temperature \\
$S(x)$, $x_{\rm RH}$ & \texttt{dkS[x]} & (7.3) & fraction of the freeze-in yield made below $T_{\rm RH}$; $x_{\rm RH}=m_h/T_{\rm RH}$ \\
$K_\nu$ & \texttt{BesselK[ν, ·]} &  & modified Bessel function of the second kind \\
$M_{\rm Pl}^{\rm (unreduced)}$ & 1.22091\textbackslash{}*\textasciicircum{}19 & $\sqrt{8\pi}M_{\rm Pl}$ & unreduced Planck mass (freeze-in formula) \\
$N$ & \texttt{nch} &  & number of decay channels (lepton flavours, brane spinors) \\
$m_*$ & \texttt{dkMstar[N]} & (7.4) & largest mass for which freeze-in and stability are compatible \\
$Y_{\rm tower}$, $\rho_{\rm tower}/s$ & \texttt{dkTowN}, \texttt{dkTowE} & (7.5) & Kaluza–Klein tower made through the Yukawa portal \\
$I_{hp}$, $I_{hp}^E$ & \texttt{dkIhp}, \texttt{dkIhpE} & (7.6) & phase-space integrals of the Higgs-portal pairs \\
$\lambda(1,x_1,x_2)$ & — & (7.6) & Källén function \\
$T_V$ & \texttt{dkTV} & $(M_8^{12}/M_{\rm Pl})^{1/11}$ & temperature above which pairs of gravitational strength are made efficiently \\*
\end{longtable}}
\begin{align*}
O&=\frac f{1-f}\,\frac{T_{\rm RH}}{\frac43mY},\qquad \frac{\rho_{\rm DM}}{\rho_{\rm rad}}=\frac43\,\frac{mY}{T},&&(7.1) \\
\Gamma(\phi\to f\bar f)&=\frac{N_cm_f^2m_r}{24\pi M_{\rm Pl}^2}\Big(1-\frac{4m_f^2}{m_r^2}\Big)^{3/2},\quad \Gamma(\phi\to gg)=\frac{b_{\rm QCD}^2\alpha_s^2m_r^3}{96\pi^3M_{\rm Pl}^2},\quad \Gamma_n=\frac{3\mu_n^2m_r}{8\pi M_{\rm Pl}^2}\Big(1-\frac{4\mu_n^2}{m_r^2}\Big)^{3/2},&&(7.2) \\
S(x)&=\frac2{3\pi}\int_x^\infty t^3K_1(t)\,dt,\qquad Y=\frac{135\,g_B}{8\pi^3\cdot1.66\,g_*^{3/2}}\,\frac{M_{\rm Pl}^{\rm (unreduced)}\Gamma}{m_B^2}\,S(x_{\rm RH}),&&(7.3) \\
m_*&=\Big(\frac{Y(c_y=1)\,K}{mY\,t_0}\Big)^{1/2},&&(7.4) \\
\sum_n\Big(1-\frac{\mu_n^2}{m_h^2}\Big)^2&=\frac{V_\perp m_h^4}{192\pi^2},\qquad \sum_n\mu_n\Big(1-\frac{\mu_n^2}{m_h^2}\Big)^2=\frac{V_\perp m_h^5}{315\pi^2},&&(7.5) \\
I_{hp}&=\iint_{u_1+u_2<1}u_1^3u_2^3\big(1-(u_1+u_2)^2\big)^{3/2}\big(1-(u_1-u_2)^2\big)^{1/2}du_1du_2,&&(7.6) \\
\lambda(1,x_1,x_2)&=(1-x_1-x_2)^2-4x_1x_2 .
\end{align*}
\Needspace{12\baselineskip}\section*{8. Dark energy}\addcontentsline{toc}{section}{8. Dark energy}
{\small
\begin{longtable}{>{\raggedright\arraybackslash}p{\dimexpr 0.1700\linewidth-2\tabcolsep\relax}>{\raggedright\arraybackslash}p{\dimexpr 0.2100\linewidth-2\tabcolsep\relax}>{\raggedright\arraybackslash}p{\dimexpr 0.3400\linewidth-2\tabcolsep\relax}>{\raggedright\arraybackslash}p{\dimexpr 0.2800\linewidth-2\tabcolsep\relax}}
\toprule
\textbf{symbol} & \textbf{Mathematica} & \textbf{definition} & \textbf{meaning} \\
\midrule\endfirsthead
\toprule
\textbf{symbol} & \textbf{Mathematica} & \textbf{definition} & \textbf{meaning} \\
\midrule\endhead
\endfoot
\bottomrule\endlastfoot
$\rho$, $\rho_S$, $k_F$ & \texttt{rho7[kF]}, \texttt{c7 iS7} & (8.1) & density, scalar density, Fermi momentum of uniform matter (7 space dimensions) \\
$m^*$, $g_S$, $g_V$ & \texttt{gap7}, \texttt{gS0}, \texttt{gV0} & (8.1) & effective mass; scalar and vector couplings ($-10^5$, $10^4$) \\
$\varepsilon$, $\varepsilon_F$, $\mu$ & \texttt{eps7}, \texttt{fermi7} & (8.1) & energy density, Fermi energy, chemical potential \\
$p$, $w$ & \texttt{dkPres[kF]}, \texttt{dkWtab} & (8.2) & pressure and equation of state \\
$k_{F0}$, $\rho_0$ & \texttt{kF0} & $\varepsilon/\rho$ minimal & saturation \\
$\rho_\Lambda$ & \texttt{dkRhoL} & $3H_0^2M_{\rm Pl}^2\Omega_\Lambda$ & dark-energy density today \\
$\rho_{\rm Cas}$ & \texttt{dkCas[m\texorpdfstring{\ensuremath{\ell}}{ℓ}]} & (8.3) & Casimir (vacuum) energy per brane volume on the torus \\
$Z$ & \texttt{dkZ48} & (8.3) & Epstein zeta value $\sum_{\vec n\ne0}|\vec n|^{-8}$ \\
$\zeta(s)$ & \texttt{Zeta} & $\sum_{k\ge1}k^{-s}$ & Riemann zeta function \\*
\end{longtable}}
\begin{align*}
\rho&=\frac{k_F^7}{105\pi^4},\quad c_7=\frac1{15\pi^4},\quad \rho_S=c_7\int_0^{k_F}\frac{k^6m^*dk}{\sqrt{k^2+m^{*2}}},\quad m^*=m+g_S\rho_S,&&(8.1) \\
\varepsilon&=c_7\int_0^{k_F}k^6\sqrt{k^2+m^{*2}}dk-\frac{(m^*-m)^2}{2g_S}+\frac{g_V}2\rho^2,\qquad \varepsilon_F=g_V\rho+\sqrt{k_F^2+m^{*2}}, \\
p&=\rho^2\frac{d}{d\rho}\Big(\frac\varepsilon\rho\Big)=\mu\rho-\varepsilon,\qquad \mu=\frac{d\varepsilon}{d\rho}=\varepsilon_F,\qquad w=\frac p\varepsilon,&&(8.2) \\
\rho_{\rm Cas}^{(1)}&=-\frac2{\ell_\perp^4}\sum_{\vec n\ne0}\Big(\frac{m\ell_\perp}{2\pi|\vec n|}\Big)^4K_4(m\ell_\perp|\vec n|)\xrightarrow{m\to0}-\frac{\Gamma(4)Z}{\pi^4\ell_\perp^4},&&(8.3) \\
Z&=\sum_{k\ge1}\frac{r_4(k)}{k^4}=8(1-4^{-3})\zeta(4)\zeta(3).
\end{align*}
$\rho_{\rm Cas}^{(1)}$ is per real bosonic degree of freedom with periodic boundary conditions; a fermion contributes with the opposite sign; dirac16complex has 16 real degrees of freedom on shell ($\rho_{\rm Cas}=-16\rho^{(1)}_{\rm Cas}$), the eight-dimensional graviton 20.

\Needspace{12\baselineskip}\section*{9. Symbols that look alike}\addcontentsline{toc}{section}{9. Symbols that look alike}
\begin{itemize}
\item $W$: the world brane $\lbrace y=0\rbrace$, and $W(a,b)$, the comoving energy of a source (always written with its arguments).
\item $h$: the induced metric $h_{\mu\nu}$ of the brane, the reduced Hubble constant $h=0.674$, and the Higgs boson $h$ (mass $m_h$).
\item $\mu$: the reduced mass $mm_N/(m+m_N)$ (section 6), the chemical potential (section 8), and with an index the Kaluza–Klein mass $\mu_n$ and the portal mass $\mu_P$.
\item $\lambda$: the Higgs-portal coupling $\lambda_P$ and $\lambda_{\rm eff}$, the Källén function $\lambda(1,x_1,x_2)$; the Higgs self-coupling is not used here.
\item $N$: the number of decay channels; $N(\mu)$ the number of Kaluza–Klein levels lighter than $\mu$; $N_c$ the number of colours.
\item $S$: the Boltzmann factor $S(x)$ (section 7), and $S_P$, the portal term of the action, (2.5).
\item $K$: the lifetime constant $K$ (section 6); $K_\nu$ the Bessel function; $K_U$ the stiffness of $U(b)$; $K$ as a multi-index in $\Theta^K$.
\item $\sigma$: the cross sections $\sigma_{\rm SI}$, $\sigma$; the divisor sum $\sigma(k)$; $\sigma_r$ the sign exponent in $\Theta^K$; $\sigma_2$ the Pauli matrix.
\item $g$: the number of internal states (Pauli bound), $g_*$, $g_{*s}$, $g_B$ (degrees of freedom), $g_r$ (self-couplings), $g_S$, $g_V$ (Kohn–Sham couplings), $g_n$ (radion couplings of the levels), $g_{\mu\nu}$ (metric).
\item $\Gamma$: decay rates; $\Gamma_A$ the Dirac matrices; $\Gamma(4)=6$ the Gamma function.
\end{itemize}
\Needspace{12\baselineskip}\section*{10. Mathematica names of section 24}\addcontentsline{toc}{section}{10. Mathematica names of section 24}
{\footnotesize
\begin{longtable}{>{\raggedright\arraybackslash}p{\dimexpr 0.4500\linewidth-2\tabcolsep\relax}>{\raggedright\arraybackslash}p{\dimexpr 0.5500\linewidth-2\tabcolsep\relax}}
\toprule
\textbf{name} & \textbf{meaning} \\
\midrule\endfirsthead
\toprule
\textbf{name} & \textbf{meaning} \\
\midrule\endhead
\endfoot
\bottomrule\endlastfoot
\texttt{dkOch2}, \texttt{dkh}, \texttt{dkOc}, \texttt{dkHbarC} & $\Omega_ch^2=0.120$, $h=0.674$, $\Omega_c$, $\hbar c$ in GeV cm \\
\texttt{dkRhoc}, \texttt{dkGs0}, \texttt{dks0}, \texttt{dkYm} & $\rho_c/h^2$, $g_{*s0}$, $s_0$, $mY$ \\
\texttt{dkMD}, \texttt{dkPauli} & $M_D$; Pauli bound for $g=16$ (keV) \\
\texttt{dkSt}, \texttt{mDM}, \texttt{kaP}, \texttt{kbP} & stress tensor of a massive particle; its mass and comoving momenta (symbols) \\
\texttt{dkDbeta} & $\delta\beta$ \\
\texttt{dkLam}, \texttt{dkYeff}, \texttt{dkGr}, \texttt{dkGeV2g} & $\lambda_{\rm eff}$, $y_{\rm eff}$, $G_r$ per unit coupling; grams per GeV \\
\texttt{dkSigSI}, \texttt{dkSigM}, \texttt{dkCrMax} & $\sigma_{\rm SI}$, $\sigma/m$, largest $c_r$ allowed by the Bullet Cluster \\
\texttt{dkMuP}, \texttt{dkTauG}, \texttt{dkTau3}, \texttt{dkAge}, \texttt{dkMmax} & $\mu_P$, lifetimes, $t_0$ in s, $m_{\max}$ \\
\texttt{dkF}, \texttt{dkOver}, \texttt{dkFracBulk} & $f$, $O(T_{\rm RH})$, $\Gamma_{\rm bulk}/(\Gamma_{\rm bulk}+\Gamma_{\rm brane})$ \\
\texttt{dkTstar}, \texttt{dkS}, \texttt{dkSstar}, \texttt{dkC} & $T_*$, $S(x)$, $S(m_h/T_*)$, the freeze-in prefactor \\
\texttt{dkYfi}, \texttt{dkYhp}, \texttt{dkK}, \texttt{dkMstar}, \texttt{dkCyNeed}, \texttt{dkCyMax} & zero-mode yields (Yukawa, Higgs portal), $K$, $m_*$, $c_y$ needed and allowed \\
\texttt{dkTowN}, \texttt{dkTowE}, \texttt{dkLat}, \texttt{dkLatE}, \texttt{dkCyTow} & tower sums (continuum), lattice sums ($R=100$), $c_y$ bound \\
\texttt{dkIhp}, \texttt{dkIhpE}, \texttt{dkPairs}, \texttt{dkRhp}, \texttt{dkCpNeed} & $I_{hp}$, $I^E_{hp}$, $(V_\perp m_h^4/(8\pi^2))^2$, mass per entropy of the pairs, $c_P$ needed \\
\texttt{dkTV} & $T_V$ \\
\texttt{dkPE}, \texttt{dkPres}, \texttt{dkWtab}, \texttt{dkWmin}, \texttt{dkKmin} & $(p,\varepsilon)$, $p$, table of $(k_F,p,w)$, $w_{\min}$ and its $k_F$ \\
\texttt{dkZ48}, \texttt{dkLp}, \texttt{dkRhoL}, \texttt{dkC0}, \texttt{dkCas}, \texttt{dkXL} & $Z$, $\ell_\perp$, $\rho_\Lambda$, $\Gamma(4)Z/\pi^4$, $\rho_{\rm Cas}$, $m\ell_\perp$ at $\rho_{\rm Cas}=\rho_\Lambda$ \\*
\end{longtable}}
\end{document}
